% SOURCE AUTHORITY:
% expose_VII_body.tex lines 873--899; scan indices 181--182; source folios 170--171.
% The conventional inverse-image notation f^* is used for the scan's typewritten f*.
% Page-header tokens at source lines 889 and 904 are excluded as running matter.
% Hard stop before source line 907, section 2.9.

2.8. \underline{Change of topos. Localization.}

Let
\[
 f:\underline T'\longrightarrow\underline T
\]
be a morphism of topoi. The inverse-image functor $f^*$ transforms a
biextension of $(P,Q)$ by $G$ into a biextension $f^*(E)$ of
$(f^*P,f^*Q)$ by $f^*G$, as always for a structure expressible in
terms of finite projective limits and arbitrary inductive limits. To
prove compatibility of the two extension structures on $f^*(E)$
induced by those on $E$, one expresses compatibility by commutativity
of (2.1.1) in the universal case
\[
 S=P\times P\times Q\times Q,
\]
where $p,p',q,q'$ are the four projections. As expected,
$E\mapsto f^*(E)$ gives a canonical functor
\[
 f^*:\operatorname{BIEXT}(\underline T)
 \longrightarrow
 \operatorname{BIEXT}(\underline T').
\]

In what follows we shall use without further mention the fact that
every construction made with biextensions is compatible with inverse
images under morphisms of topoi. This applies in particular to the
product $EE'$ and inverse $E^{-1}$ considered in 2.5.

One may take for $\underline T'$ an \underline{induced topos}
$\underline T_{/S'}$, and for $f$ the canonical morphism, called the
localization morphism (SGA 4 IV 5),
\[
 f=i_{S'}:\underline T_{/S'}\longrightarrow\underline T.
\]
Then $f^*(E)$ is called the \underline{biextension induced} by $E$ on
$S'$, and is also denoted, as usual, by $E|_{S'}$ or $E_{S'}$.

Since restriction has the usual transitivity properties, the
categories $\operatorname{BIEXT}(\underline T_{/S})$, for variable
$S\in\underline T$, form the fibres of a \underline{fibred category}
over $\underline T$ (SGA 1 VI 8). In fact, this category is even a
\underline{stack} over $\underline T$ in the sense of Giraud [2]:
morphisms and objects in its fibres can be glued. For example, to give
a biextension over $S$ up to canonical isomorphism, it suffices to give
a covering family (that is, an epimorphic family) $S_i\longrightarrow
S$, a biextension $E_i$ over each $S_i$ (an object of
$\underline T_{/S_i}$), and, for every pair $(i,j)$, an isomorphism
between $E_i|_{S_j}$ and $E_j|_{S_i}$, subject to the usual cocycle
condition.

We call this the \underline{stack of biextensions over $\underline T$}.
Replacing $\underline T$ in (2.3.3), and hence also in (2.3.4) and
(2.3.5), by $\underline T_{/S}$ for variable $S$ gives morphisms of
stacks over $\underline T$. The same applies to the operations $E.E'$
and $E^{-1}$ of (2.5.1) in the categories
$\operatorname{BIEXT}(P_S,Q_S;G_S)$ for variable $S$, where $P,Q,G$
are fixed commutative groups over $\underline T$.

